\section{Withdrawal Contracts and Deferred Work}
\label{sec:background}

\subsection{Four contracts, not one}
\label{sec:contracts}

``The user said stop'' covers several different promises. We separate four, following the distinctions that consent
studies~\cite{nguyen2021sharefirst,nguyen2022freely,kollnig2021faitaccompli} and data-deletion work~\cite{reardon2013deletion,garg2020deletion,shastri2020gdprbench}
already draw, and evaluate every fixture against the contract it states (Table~\ref{tab:contracts}).

\begin{table}[t]
\centering
\caption{Withdrawal contracts. ``Pending'' says whether data collected \emph{before} the withdrawal is covered.}
\label{tab:contracts}
\mtab
\begin{tabular}{@{}>{\raggedright\arraybackslash}p{2.2cm}>{\raggedright\arraybackslash}p{1.6cm}>{\raggedright\arraybackslash}p{1.15cm}>{\raggedright\arraybackslash}p{2.8cm}@{}}
\toprule
\mhead Contract & Purposes revoked & Pending & Violation evidence \\
\midrule
OS permission revocation & none (new access only) & no & none from prior data; residue is \emph{not} a violation \\
In-app opt-out (stop collecting) & all & stated & new collection; prior data only if the app says ``discard'' \\
Processing restriction & named purpose & yes & use for that purpose; other purposes stay legal \\
Deletion request & all & yes & any later egress of the data or its derivatives \\
\bottomrule
\end{tabular}
\end{table}

An \emph{OS permission revocation} (\texttt{pm revoke}, or the settings toggle) only stops future access to a source.
Android kills the process, but nothing in the platform contract obliges the app to forget what it already read,
so a later upload of earlier location fixes is residue, not a violation~\cite{felt2011demystified,wang2022aper}.
An \emph{in-app opt-out} is the app's own promise. Some wordings stop collection only (``we will stop collecting''); others
also discard what is queued (``we will no longer send usage data''). A \emph{processing restriction} withdraws one purpose,
such as personalisation, while storage and backup stay permitted. A \emph{deletion request} covers the data and anything derived
from it. Throughout, a submission is a \emph{prohibited submission} only if its contents contradict the contract the app states, at the time
the bytes leave the device. Stored residue alone never counts. We make no legal claims.

\subsection{Deferred work survives the withdrawal}
\label{sec:deferred}

Persisted \texttt{JobScheduler} jobs are written to
the system's job store and restored after reboot~\cite{android_jobscheduler}. WorkManager keeps its own SQLite database and reschedules work
after app restarts and reboots~\cite{android_workmanager}. Alarms do not survive reboot, so apps re-arm them from their own state. SyncAdapters
run on the system's schedule. Analytics and crash-reporting SDKs add another layer: they keep an outbox, and many flush it when the app
starts.

On the Android~13 image we use (\S\ref{sec:setup}), we observed the following lifecycle facts directly over \lcTrials{} campaign trials
(Table~\ref{tab:lifecycle}). Each one shapes which enforcement points can see a queued upload.
\begin{table}[t]
\centering
\caption{Measured lifecycle behavior of a persisted upload job (\vname{none}, AVD~A). Counts of trials; ``by-value job'' is T2's summary job, whose data exists only in its extras.}
\label{tab:lifecycle}
\mtab
\begin{tabular}{@{}lrrrrr@{}}
\toprule
\mhead Schedule & trials & job ran & re-enqueued & by-value job ran & median s \\
\midrule
\textsf{pen} & 50 & 50 & 0 & 50 & 43 \\
\textsf{kill} & 50 & 50 & 0 & 50 & 43 \\
\textsf{fstop} & 50 & 0 & 50 & 0 & 19 \\
\textsf{reb} & 50 & 50 & 0 & 50 & 51 \\
\textsf{rev} & 50 & 50 & 0 & 50 & 24 \\
\bottomrule
\end{tabular}
\end{table}

\begin{itemize}[leftmargin=*,nosep]
\item \textbf{Reboot.} Persisted jobs are restored and run after boot. A check made at enqueue time is therefore arbitrarily old when the job runs.
\item \textbf{Process death} (\texttt{kill -9}, or a runtime-permission revocation, which kills the app). Jobs stay scheduled, and JobScheduler restarts the process to run them.
\item \textbf{Force-stop} removes the app's scheduled jobs. When the user next opens the app, an outbox-driven startup path re-enqueues the work. A job whose \emph{only} copy of the data was in its extras is lost.
\item \textbf{Network loss mid-job.} An upload loop retries the batch it already prepared. A check made at job start ran before the withdrawal.
\end{itemize}
Withdrawal is a user action at an arbitrary time, so it can land in any of these windows. Our question is which enforcement point
still sees it.
